

Or si @@M0@@ est un @@M1@@-anneau, le diagramme :
@@M2@@
est commutatif d'après 7.3, puisque @@M3@@ est un @@M4@@-homomorphisme. On en déduit aussitôt les relations cherchées.

ii) Si @@M5@@ est de caractéristique nulle, l'application « dérivée logarithmique » est un isomorphisme de @@M6@@ sur @@M7@@, car on peut intégrer. Par suite, deux éléments de @@M8@@ sont égaux si, et seulement si pour tout @@M9@@, ils ont même image par @@M10@@.

Pour montrer que @@M11@@ est un homomorphisme d'anneaux, il faut montrer la commutativité du carré de gauche du diagramme (7.5.2), et d'après la remarque précédente, il suffit de montrer que pour tout @@M12@@, le diagramme (7.5.2) est commutatif, puisque le carré de droite l'est toujours. Or ceci n'est autre que les relations @@M13@@.

Pour montrer que @@M14@@ est un @@M15@@-homomorphisme, il suffit de montrer que le diagramme (7.5.4) est commutatif pour tout @@M16@@, d'après 7.3. Or, d'après la remarque qui précède, cela équivaut à montrer pour tous @@M17@@, que @@M18@@, c'est-à-dire les relations @@M19@@.

Enfin, la relation @@M20@@ s'obtient par intégration à partir des relations @@M21@@.

\subsection*{BIBLIOGRAPHIE}
\begin{enumerate}[label={[\arabic*]}]
\item A. Grothendieck, La théorie des classes de Chern, Bull. Soc. math. France, 86, 1958, p. 137 à 154.
\item F. Hirzebruch, Neue topologische Methoden in der algebraischen Geometrie.
\end{enumerate}



% Exposé VI projective-bundle macro additions
\providecommand{\Sym}{\operatorname{Sym}}
\providecommand{\Proj}{\operatorname{Proj}}
\providecommand{\Projz}{\operatorname{Proj}_0}
\providecommand{\Spec}{\operatorname{Spec}}
\providecommand{\Comp}{\operatorname{Comp}}
\providecommand{\Dplus}{D_{+}}
\renewcommand{\id}{\operatorname{id}}
\providecommand{\cG}{\mathcal G}
\providecommand{\cA}{\mathcal A}
\providecommand{\cB}{\mathcal B}
\providecommand{\cL}{\mathscr L}
\providecommand{\cM}{\mathscr M}
\providecommand{\cN}{\mathscr N}
\providecommand{\cE}{\mathscr E}
\providecommand{\cK}{\mathscr K}
\providecommand{\cP}{\mathscr P}
\providecommand{\cQ}{\mathscr Q}
\providecommand{\OO}{\mathcal O}
\providecommand{\rk}{\operatorname{rk}}
\providecommand{\colim}{\operatorname*{colim}}



